% =====================================================================================================
\section{Result 1: the posterior finds the true model}\label{sec:model-result}

\inwords{If the true model is small and its coefficients are fixed and nonzero, and no small group of predictors is
collinear, then the posterior puts all but a fraction of order $\rn$ of its mass on the true model. The highest posterior
model and the median probability model then equal the true model with probability tending to one, and every inclusion
probability is within order $\rn$ of $1$ for a true predictor and of $0$ for a null one. We state this first for Gaussian
rows, the case of our simulation designs, and then for a fixed design, with conditions on the sample.}

\why{It is the basic guarantee of a Bayesian selection method, here for the prior exactly as the software sets it: with
enough data the posterior points to the right predictors, and its inclusion probabilities can be read as near-certain
inclusion or exclusion.}

\begin{assumption}{FSb}{Fixed signals, raw scale}
$b_{\mathrm{lo}}\le\abs{\beta^*_j}\le b_{\mathrm{hi}}$ for every $j\in\Sz$, with fixed $0<b_{\mathrm{lo}}\le b_{\mathrm{hi}}<\infty$.
\end{assumption}

\begin{theorem}[Posterior concentration, random Gaussian design]\label{thm:post-random}
Assume \as{GL}, \as{DIM}, \as{SP}, \as{GR}, \as{FSb} and \as{EFS}. Then, under the joint law of the design and the noise,
\begin{gather}
  1-\pis{\Sz}=\OP\Big(\sqrt{\frac{\log n}{n}}\Big),\qquad \P\big(\HPM=\MPM=\Sz\big)\to1,\label{eq:post-random}\\
  \max_{j\in\set{p}}\big|\PIP_j-\ind{j\in\Sz}\big|=\OP\Big(\sqrt{\frac{\log n}{n}}\Big).\nonumber
\end{gather}
\end{theorem}

\asmlist{(1)~\as{GL} Gaussian linear model; (2)~\as{DIM} $p\le n^A$; (3)~\as{SP} fixed $\sz$ [\SA{3.1}]; (4)~\as{GR}
Gaussian rows; (5)~\as{FSb} fixed nonzero coefficients; (6)~\as{EFS} population eigenvalue floor on sets of
$O(\log n)$ predictors [\SA{3.2}].}

\status{\textbf{Proved and independently checked}: the fixed-design result it rests on (Theorem~\ref{thm:post-fixed},
in-probability form) and the transfer from population to sample conditions are each proved and independently checked.
Registry: Appendix~\ref{app:status}, row 1.2.}

No upper bound on the eigenvalues of $\Sig$ is needed, so strong common factors are allowed, and $p\gg n$ is allowed. In
primitive terms: a Gaussian linear model with i.i.d.\ Gaussian rows whose correlation matrix has smallest eigenvalue bounded below, a fixed number of
fixed nonzero coefficients and $p\le n^A$ suffice.

\paragraph{Fixed design.} The same conclusions hold for a deterministic design under conditions on the sample. They
use four more assumptions.

\begin{assumption}{BC}{Bounded coefficients}
$\max_{j}\abs{\theta^*_j}\le\bar\theta$ for a fixed $\bar\theta<\infty$. \hfill [\SA{3.2}]
\end{assumption}

\begin{assumption}{EF}{Eigenvalue floor on sets of bounded size}
$\lmin(\Rhat_S)\ge\varphi_0^2$ for every model $S$ with $\abs{S}\le2(A+1)\sz+2$, for a fixed $\varphi_0\in(0,1]$. \hfill [\SA{3.2}]
\end{assumption}

\begin{assumption}{NC}{No near-exact collinearity among sets of logarithmic size}
$\lmin(G_S)\ge g_0$ for every model $S$ with $\abs{S}\le3(A+1)(\sz+1)\log n$, for a fixed $g_0>0$; on the correlation
scale, a floor $g_0/(n-1)$ that vanishes as $n$ grows. \hfill [\SA{3.2}]
\end{assumption}

\begin{assumption}{BM}{Beta-min}
Either (a) $\tmin\ge C_\beta\,\sigs\rn$ with a fixed $C_\beta>0$ such that $C_\beta^2\varphi_0^2>c_0(\sz,A)$, an explicit
constant (Appendix~\ref{app:constants}); or (b) $\tmin\sqrt{n/\log n}\to\infty$. \hfill [\SA{3.2}]
\end{assumption}

\begin{theorem}[Posterior concentration, fixed design]\label{thm:post-fixed}
Let the design be fixed and assume \as{GL}, \as{DIM}, \as{SP}, \as{BC}, \as{EF}, \as{NC} and \as{BM}. Then, along every
sequence that satisfies these assumptions,
\begin{equation}\label{eq:post-op}
  \max_{j\in\set{p}}\big|\PIP_j-\ind{j\in\Sz}\big|\ \le\ 1-\pis{\Sz}\ =\ \OP\Big(\sqrt{\frac{\log n}{n}}\Big),\qquad
  \P\big(\HPM=\MPM=\Sz\big)\to1 .
\end{equation}
Moreover, there are $n_0$ and $C_{\mathrm E}$, depending only on $(\sz,A,\sigs^2,\bar\theta,\varphi_0,g_0,C_\beta)$, such that for
every $n\ge n_0$ and every design and $\tstar$ that satisfy the assumptions at that $n$,
\begin{equation}\label{eq:post-exp}
  \E\big[1-\pis{\Sz}\big]\ \le\ C_{\mathrm E}\sqrt{\frac{\log n}{n}},\qquad \P\big(\HPM=\MPM=\Sz\big)\ \ge\ 1-4C_{\mathrm E}\sqrt{\frac{\log n}{n}} .
\end{equation}
The expectation is over the noise, with the design fixed. Under form (b) of \as{BM}, $n_0$ also depends on how fast
$\tmin\sqrt{n/\log n}$ diverges.
\end{theorem}

\asmlist{(1)--(3)~\as{GL}, \as{DIM}, \as{SP} [\SA{3.1}]; (4)~\as{BC} bounded coefficients; (5)~\as{EF} eigenvalue floor
on sets of bounded size; (6)~\as{NC} no near-exact collinearity; (7)~\as{BM} beta-min [\SA{3.2}].}

\status{The in-probability conclusions~\eqref{eq:post-op}: \textbf{proved and independently checked}. The expectation
bound~\eqref{eq:post-exp}: \textbf{proved and checked, conditional on one supporting result} (a proposition giving a lower
bound on the posterior miss) \textbf{that is proved but not yet independently checked}; its other input, a lemma bounding
the tail of the posterior model size, is now independently checked. Registry: Appendix~\ref{app:status}, row 1.1.}

\subsection{Result 5: the rate cannot be improved}\label{sec:sharp}

\inwords{On every design, no model of at most $\min(K-1,p/2)$ predictors has posterior mass closer to one than about
$(\abs{S}+1)/(3\sqrt n)$; on two families of designs the expected posterior miss is at least of order $\rn$. At simulated sizes the posterior spreads over the true model and
its one-predictor extensions.}

\why{It shows that the rate of Result~1 is the right one, and why the posterior is not a point mass at simulated sizes
such as $n=200$.}

\begin{remark}[The rate cannot be improved, and what it means at moderate $n$]\label{rem:sharp}
\textup{(a)} \emph{Every design.} For every $n\ge3$, every standardised design with $p\ge2$, every response and every
model $S$ with $\abs{S}\le\min(K-1,p/2)$,
\begin{equation}\label{eq:sharp-a}
  \pis{S}\ \le\ \Big(1+\frac{\abs{S}+1}{3\sqrt n}\Big)^{-1}.
\end{equation}
So no model of that size, the true one included, can carry posterior mass closer to one than about
$(\abs{S}+1)/(3\sqrt n)$. For true models with $\sz\le\min(K-1,p/2)$, no design gives a rate faster than
$n^{-1/2}$, and a posterior miss of order $\rn$ is possible only if $\sz=O(\sqrt{\log n})$. The reason is that the prior and the Occam factor keep every one-predictor extension of a model
alive.

\textup{(b)} \emph{Witness designs.} Assume \as{GL}, that $\sz$ is fixed, the energy bounds $\norm{\mu}^2\le nM_\mu$ and
$\norm{\tstar}^2\le nM_\theta$ for fixed $M_\mu$, $M_\theta$, and that the null predictors are either orthogonal (to each
other and to the true ones) or have i.i.d.\ $\Normal(0,1)$ raw entries independent of the true predictors and of the noise, with $p\le n^A$.
Then, for every $p>\sz$, bounded or growing,
where $\pi=3.14\ldots$ (the number, not the posterior) in the next display,
\begin{equation}\label{eq:sharp-b}
  \E\big[1-\pis{\Sz}\big]\ \ge\ (1-o(1))\,(\sz+1)\,\frac{2(p-\sz)}{2p-1-\sz}\,
  \sqrt{\frac{\log\!\big(p\sqrt n/(\sz+1)\big)}{\pi n}},
\end{equation}
and, if moreover $p\to\infty$, $1-\pis{\Sz}\ge(1-\oP(1))\,(\sz+1)\sqrt{\log p/(\pi n)}$. On fully
orthogonal designs ($\sz<p\le n-1$) with fixed nonzero coefficients, the expectation bound~\eqref{eq:post-exp} applies
too, so there $\E[1-\pis{\Sz}]$ is of exact order $\rn$.

\textup{(c)} \emph{At moderate $n$} (a heuristic, not a theorem). The posterior miss is carried mainly by the models
$\Sz\cup\{j\}$ that add one null predictor. Each has posterior odds against $\Sz$ of about
$\frac{\sz+1}{2p}\,n^{-1/2}e^{T_j^2/2}$ (prior, Occam factor and fit), with $T_j^2$ approximately $\chi^2_1$. The factor
$e^{T_j^2/2}$ has no finite mean; over $p$ null predictors the largest $T_j^2$ is about $2\log p$, and the $p$ factors
average about $2\sqrt{\log p/\pi}$. Together these models therefore carry about $(\sz+1)\sqrt{\log p/(\pi n)}$ times the
weight of $\Sz$: the expression of~(b), here as a ratio to the weight of $\Sz$ rather than a bound on the posterior miss. With $\sz=5$, $n=200$ and $p=10^4$ this is about $0.7$. In single
synthetic draws (AR(0.6), $\sz=5$) the measured total weight of these models, relative to that of $\Sz$, was $1.0$--$1.2$
at $n=200$ and $p=10^4$--$10^6$, and $0.52$ at $n=400$, $p=10^4$. So at simulated sizes the posterior is spread over $\Sz$
and many one-predictor extensions of it.
\end{remark}

\asmlist{(a): none. (b): \as{GL}, a fixed $\sz$ [\SA{3.1}], and the witness designs, which describe where the rate is
attained and are not conditions of consistency.}

\status{(a) and (b): \textbf{proved and independently checked}; the exact-order statement after (b) uses the
expectation bound~\eqref{eq:post-exp} and has its status. Registry: Appendix~\ref{app:status},
row 1.3. (c): a heuristic computed by the Theory department, with one-draw measurements from the engineering
experiments; no registry row.}
